\PDFsubsection{SCFT of L--S TT Copolymers}

    \begin{figure}[!htb]
        \centering
        \includegraphics[]{./Figure/Illustrate.png}
        \caption{
            (a) Illustration of an L--S TT copolymer, in which a sliding side chain is grafted onto the A block of the main chain at the sliding point through a rotaxane structure.
            The black dashed line denotes the chain contour coordinate, and the red and blue regions represent the A and B components, respectively.
            (b--d) Topological states at different $f$: linear AB diblock at $f=0$, intermediate A$_2$AB topology for $0<f<1$, and three-arm A$_2$B star topology at $f=1$.
        }
        \label{fig:Architecture}
    \end{figure}

    We consider a melt of $n$ L--S TT block copolymers in a three-dimensional volume $V$.
    As shown in Fig.~\ref{fig:Architecture}a, each copolymer contains a total number of segments
    $N = N_{\mathrm{A1}} + N_{\mathrm{A2}} + N_{\mathrm{A3}} + N_{\mathrm{B}}$.
    Here, A1 denotes the side chain composed of A segments, which is grafted onto the main chain through a rotaxane unit and can freely slide along the A block of the main chain; its number of segments is $N_{\mathrm{A1}}$.
    The main chain consists of an A block and a B block, and the position of the rotaxane sliding point further divides the A block into two sub-blocks, A2 and A3.
    The A block between the free end of the main chain and the grafting point is referred to as the tail, denoted by A2, with $N_{\mathrm{A2}}$ segments.
    The A block between the grafting point and the junction with the B block is referred to as the trunk, denoted by A3, with $N_{\mathrm{A3}}$ segments.
    The B block is denoted by B and contains $N_{\mathrm{B}}$ segments.
    Therefore, the total number of A segments on the main chain is $N_{\mathrm{A2}} + N_{\mathrm{A3}}$.
    The overall volume fraction of the A component is defined as
    $\bar{\phi}_{\mathrm{A}} = 1 - \bar{\phi}_{\mathrm{B}}$, where
    $\bar{\phi}_{\mathrm{B}} = N_{\mathrm{B}}/N$ is the volume fraction of the B component.
    The side-chain fraction is defined as
    $\phi_{\mathrm a} = N_{\mathrm{A1}}/(N_{\mathrm{A1}} + N_{\mathrm{A2}} + N_{\mathrm{A3}})$,
    which characterizes the fraction of A segments contained in the side chain.
    To describe the motion of the sliding point along the main chain, we introduce a sliding variable $\kappa \in [0, N_{\mathrm{A2}} + N_{\mathrm{A3}}]$, which denotes the number of segments from the free end of the main chain to the grafting point.
    The corresponding dimensionless sliding degree of freedom is defined as
    \begin{equation}
        f = \frac{\kappa}{N_{\mathrm{A2}} + N_{\mathrm{A3}}},
    \end{equation}
    where $N_{\mathrm{A2}} + N_{\mathrm{A3}}$ is the total length of the slidable A block on the main chain.
    As $f$ varies from 0 to 1, the copolymer can continuously transform among the AB linear topology (Fig.~\ref{fig:Architecture}b), the A$_2$AB intermediate star-like topology (Fig.~\ref{fig:Architecture}c), and the A$_2$B star topology (Fig.~\ref{fig:Architecture}d).
    Accordingly, the physical interpretation of $f$ as a sliding-topology variable applies only for $0<\phi_{\mathrm a}<1$; at $\phi_{\mathrm a}=0$ and $1$, the corresponding sliding branch vanishes, and the architecture reduces to a linear AB diblock.
    

    
    To describe the connectivity between different blocks, we introduce a reference contour coordinate $s\in[0,1]$ based on the limiting topology at $f=0$, following the black dashed line in Fig.~\ref{fig:Architecture}a.
    In this parametrization, $s$ is assigned to increase from the free end of the A1 side chain to the free end of the B block.
    For notational convenience in constructing the propagators, each block $m\in\{\mathrm{A1},\mathrm{A2},\mathrm{A3},\mathrm{B}\}$ is further described by a block-local contour variable $s_m$, measured along the increasing direction of the reference coordinate $s$, as indicated by the green markers in Fig.~\ref{fig:Architecture}a.
    The corresponding backward contour variable $s_m^\dagger$ is measured from the opposite end of the same block, along the decreasing direction of $s$, with $s_m+s_m^\dagger=\ell_m(f)$.
    Thus, $s_m,s_m^\dagger\in[0,\ell_m(f)]$, where $\ell_m(f)$ denotes the normalized contour length of block $m$, with $\ell_{\mathrm{A1}}=\phi_{\mathrm a}\bar{\phi}_{\mathrm A}$, $\ell_{\mathrm{A2}}(f)=f(1-\phi_{\mathrm a})\bar{\phi}_{\mathrm A}$, $\ell_{\mathrm{A3}}(f)=(1-f)(1-\phi_{\mathrm a})\bar{\phi}_{\mathrm A}$, and $\ell_{\mathrm{B}}=\bar{\phi}_{\mathrm B}$.
        

    L--S TT copolymers possess a sliding degree of freedom, and the topology $f$ of each chain is spontaneously adjusted under thermal fluctuations and intermolecular interactions.
    We denote the instantaneous topological configuration of the $i$-th chain by $f_i$.
    An integration over $f_i$ is therefore introduced into the partition function~\cite{Zhao2024Phase}:
    \begin{equation}
        \begin{split}
            Z = & \prod_{i=1}^n \int \mathcal{D}[\bm{R}_i(s)] \int_0^1 \mathrm{d}f_i\, \bigg\{ e^{-\beta(\mathcal{H}_0 + \mathcal{H}_1)} 
            \times \delta[\hat{\phi}_{\mathrm{A}}(\bm{r}) + \hat{\phi}_{\mathrm{B}}(\bm{r}) - 1]\\
            & \times \delta[\bm{R}_i(s=\phi_{\mathrm{a}}\bar{\phi}_{\mathrm{A}}+f_i(1-\phi_{\mathrm{a}})\bar{\phi}_{\mathrm{A}}) 
            - \bm{R}_i(s=\phi_{\mathrm{a}}\bar{\phi}_{\mathrm{A}})] \bigg\}.
        \end{split}
    \end{equation}
    Here, no additional intrinsic sliding potential is assigned to $f_i$, corresponding to the idealized limit in which the bare statistical weight along the sliding coordinate is uniform.
    
    The parameter $\beta$ is defined as $(k_{\mathrm B}T)^{-1}$, where $k_{\mathrm B}$ is the Boltzmann constant and $T$ is the temperature.
    The vector $\bm{R}_{i}(s)$ denotes the position of the $i$-th chain at the contour variable $s$.
    The first $\delta$ function enforces the incompressibility of the system, and the second $\delta$ function constrains one end of the sliding side chain, at $s=\phi_{\mathrm{a}}\bar{\phi}_{\mathrm{A}}$, to be connected to the main chain.
    $\mathcal{H}_0$ represents the bonded interactions of all ideal chain segments, and $\mathcal{H}_1$ represents the nonbonded interactions.
    Their specific forms are:
    \begin{equation}
        \beta \mathcal{H}_0 = 
        \frac{3}{2 N b^2} 
        \sum_{i=1}^n \int_0^{1} \mathrm{d}s 
        \left[ \frac{\mathrm{d} \bm{R}_{i}(s)}{\mathrm{d}s} \right]^2,
    \end{equation}
    \begin{equation}
        \beta \mathcal{H}_1 = 
        \chi \rho_0 \int \mathrm{d}\bm{r} \, \hat{\phi}_{\mathrm{A}}(\bm{r}) \hat{\phi}_{\mathrm{B}}(\bm{r}),
    \end{equation}
    where $b$ is the Kuhn length, $\chi$ is the Flory--Huggins interaction parameter, and $\rho_0 = nN/V$ is the average segment density.
    The quantities $\hat{\phi}_{\mathrm{A}}(\bm{r})$ and $\hat{\phi}_{\mathrm{B}}(\bm{r})$ are the local volume-fraction operators of the A and B components at position $\bm{r}$, respectively.
    They are defined by
    \begin{equation}
        \rho_0 \hat{\phi}_{\mathrm A}(\bm{r})
        \equiv
        N \sum_{i=1}^n \int_0^{\bar{\phi}_{\mathrm A}} \mathrm{d}s\,
        \delta[\bm{r} - \bm{R}_i(s)],
    \end{equation}
    \begin{equation}
        \rho_0 \hat{\phi}_{\mathrm B}(\bm{r})
        \equiv
        N \sum_{i=1}^n \int_{\bar{\phi}_{\mathrm A}}^{1} \mathrm{d}s\,
        \delta[\bm{r} - \bm{R}_i(s)].
    \end{equation}


    According to the Hubbard--Stratonovich transformation\cite{Helfand1975Theorya, Leibler1980Theory}, the partition function can be rewritten by introducing the density fields $\phi_\alpha(\bm{r})$, their conjugate fields $\omega_\alpha(\bm{r})$, and the dimensionless Lagrange-multiplier pressure field $\eta(\bm r)$ that enforces local incompressibility:
    \begin{equation}
        Z = \int \mathcal{D}[\phi_\alpha, \omega_\alpha, \eta]\, e^{-\mathcal{F}[\phi_\alpha, \omega_\alpha, \eta]}.
        \label{eq:Z}
    \end{equation}
    The corresponding dimensionless free-energy functional is given by
    \begin{equation}
        \begin{split}
            \frac{N \mathcal{F}}{\rho_0 V} 
            = &\, \frac{1}{V} \int N \chi \phi_{\mathrm{A}}(\bm{r}) \phi_{\mathrm{B}}(\bm{r})\, \mathrm{d}\bm{r} \\
            & - \frac{1}{V} \int \left[ \omega_{\mathrm{A}}(\bm{r}) \phi_{\mathrm{A}}(\bm{r}) 
            + \omega_{\mathrm{B}}(\bm{r}) \phi_{\mathrm{B}}(\bm{r}) \right] \mathrm{d}\bm{r} \\
            & + \frac{1}{V} \int \eta(\bm r)
            \left[\phi_{\mathrm A}(\bm r)+\phi_{\mathrm B}(\bm r)-1\right] \mathrm d\bm r \\
            & - \ln Q[\omega_{\mathrm{A}}, \omega_{\mathrm{B}}],
        \end{split}
        \label{eq:FreeEnergy}
    \end{equation}
    where $V$ is the system volume.
    Here, $Q$ is the single-chain partition function of the L--S TT copolymer: 
    \begin{equation}
        Q[\omega_{\mathrm{A}}, \omega_{\mathrm{B}}]
        = \int_{0}^{1} \mathrm{d} f Q_f,
        \label{eq:Partition}
    \end{equation}
    where $Q_f$ denotes the single-chain partition function for a copolymer with the grafting position fixed at a given value of $f$.


    Under the saddle-point approximation~\cite{Leibler1980Theory, Huang1998SelfConsistent, XinghuaZhang2013Effects, Zhao2024Phase}, the SCFT equations are obtained by requiring the free-energy functional $\mathcal{F}[\phi_\alpha,\omega_\alpha,\eta]$ to be stationary with respect to the density fields $\phi_\alpha(\bm r)$, their conjugate fields $\omega_\alpha(\bm r)$, and the pressure field $\eta(\bm r)$:
    \begin{equation}
        \frac{\delta \mathcal{F}}{\delta \omega_\alpha(\bm{r})} 
        = \frac{\delta \mathcal{F}}{\delta \phi_\alpha(\bm{r})} 
        = \frac{\delta \mathcal{F}}{\delta \eta(\bm{r})}
        = 0.
    \end{equation}
    This gives the self-consistent field equations:
    \begin{equation}
        \phi_\alpha(\bm{r}) = - \frac{V}{Q[\omega_\alpha]} \frac{\delta Q[\omega_\alpha]}{\delta \omega_\alpha(\bm{r})};
        \label{eq:DensityField}
    \end{equation}
    \begin{equation}
        \omega_\alpha(\bm{r}) = N \chi \, \phi_{\overline{\alpha} }(\bm{r}) + \eta(\bm r),
        \label{eq:ConjugateField}
    \end{equation}
    \begin{equation}
        \phi_{\mathrm A}(\bm r)+\phi_{\mathrm B}(\bm r)=1,
        \label{eq:Incompressibility}
    \end{equation}
    where $\alpha, \overline{\alpha} \in \{\mathrm{A}, \mathrm{B}\}$ and $\overline{\alpha} \neq \alpha$.
    In the canonical ensemble, $\eta(\bm r)$ is defined only up to an additive spatially uniform constant.
    Consistent with the convention used by \texttt{pscfpp}, this gauge freedom is removed for the periodic calculations by requiring its spatial average to vanish, $V^{-1}\int \eta(\bm r)\,\mathrm d\bm r=0$.
    

    To solve the above equations, we define the forward propagator $q_m(\bm r,s_m)$ and the backward propagator $q_m^\dagger(\bm r,s_m^\dagger)$.
    Here, $q_m(\bm r,s_m)$ represents the path integral weight for a segment to be located at $\bm r$ after propagating a contour length $s_m$ from the end of block $m$ with the smaller value of $s$.
    Similarly, $q_m^\dagger(\bm r,s_m^\dagger)$ represents the statistical weight for a segment to be located at $\bm r$ after propagating backward a contour length $s_m^\dagger$ from the other end of the block, with $s_m+s_m^\dagger=\ell_m$.
    For notational simplicity, the chain contour variable is denoted uniformly by $s$ below.
    Under this parametrization, both $q_m$ and $q_m^\dagger$ satisfy the same modified diffusion equation:
    \begin{equation}
        \frac{\partial q_m(\bm r, s)}{\partial s}
        =
        \left[
            \frac{N b^2}{6}\nabla^2 - \omega_m(\bm r)
        \right]
        q_m(\bm r, s).
    \end{equation}
    The external field $\omega_m(\bm r)$ is determined by the segment type:
    \begin{equation}
        \omega_m(\bm r)=
            \begin{cases}
                \omega_{\mathrm A}(\bm r),
                & m \in \{\mathrm{A1},\mathrm{A2},\mathrm{A3}\}, \\
                \omega_{\mathrm B}(\bm r),
                & m \in \{\mathrm{B}\}.
            \end{cases}
    \end{equation}
    The initial conditions of the propagators are determined by the block free ends and the topological connectivity between blocks.
    For the three free ends of the L--S TT copolymer, the corresponding propagators are assigned unit initial weights:
    \begin{equation}
        q_{\mathrm{A1}}(\bm r,0)=1,\quad
        q_{\mathrm{A2}}(\bm r,0)=1,\quad
        q_{\mathrm B}^\dagger(\bm r,0)=1.
    \end{equation}
    The initial conditions for the remaining blocks are determined by the chain-connectivity constraints and are given in Supporting Information Sec.~\ntxref{SI-sec:Propagator}.


    The single-chain partition function of the L--S TT copolymer can then be expressed in terms of the propagators as
    \begin{equation}
        Q_f = \frac{1}{V}
        \int \mathrm d\bm r_f \,
        q_{\mathrm{A1}}(\bm r_f,\ell_{\mathrm{A1}})
        \,q_{\mathrm{A2}}(\bm r_f,\ell_{\mathrm{A2}})
        \,q_{\mathrm{A3}}^\dagger(\bm r_f,\ell_{\mathrm{A3}}),
        \label{eq:FixedPartition}
    \end{equation}
    where $\bm r_f$ is the position of the grafting point of the $\mathrm{A_2(AB)}$ molecule in three-dimensional space.
    The first two propagators in $Q_f$ correspond to the statistical weights of the A1 side-chain and A2 tail blocks propagating from their free ends to the grafting point, respectively, whereas $q^\dagger_{\mathrm{A3}}$ represents the statistical weight of propagation from the free end of the B block, through the B block to the A/B junction, and then along the trunk block to the grafting point.
    Therefore,
    \begin{equation}
        q_{\mathrm{A3}}^\dagger(\bm r, 0) = q_{\mathrm{B}}^\dagger(\bm r, \ell_{\mathrm{B}}).
    \end{equation}


    Combining Eqs.~\eqref{eq:Partition} and \eqref{eq:DensityField}, the density of block $m$ at position $\bm r$ can be written as
    \begin{equation}
        \phi_m(\bm r)
        =
        \int_0^1
        \psi_m(\bm r;f)\,\mathrm df,
        \label{eq:BlockDensity}
    \end{equation}
    
    where $\psi_m(\bm r;f)$ denotes the topology-resolved density of block $m$ at position $\bm r$ for a given sliding configuration $f$. 
    It is obtained from the corresponding propagators as
    \begin{equation}
        \begin{aligned}
            \psi_m(\bm r;f)
            &=
            -\frac{V}{Q}
            \frac{\delta Q_f}{\delta\omega_m(\bm r)}
            \\
            &=
            \frac{1}{Q}
            \int_0^{\ell_m(f)}
            \mathrm ds_m\,
            q_m(\bm r,s_m)
            q_m^\dagger
            \left(
                \bm r,\ell_m(f)-s_m
            \right).
        \end{aligned}
        \label{eq:TopologyDensity}
    \end{equation}
    Thus, $\phi_m(\bm r)$ contains the statistical contributions from all accessible sliding configurations.
    


    For the A component, the topology-resolved density is obtained by summing the contributions from $\mathrm{A1}$, $\mathrm{A2}$, and $\mathrm{A3}$,
    \begin{equation}
        \psi_{\mathrm A}(\bm r;f)
        =
        \psi_{\mathrm{A1}}(\bm r;f)
        +
        \psi_{\mathrm{A2}}(\bm r;f)
        +
        \psi_{\mathrm{A3}}(\bm r;f),
        \label{eq:ATopologyDensity}
    \end{equation}
    and the corresponding total density is
    \begin{equation}
        \phi_{\mathrm A}(\bm r)
        =
        \phi_{\mathrm{A1}}(\bm r)
        +
        \phi_{\mathrm{A2}}(\bm r)
        +
        \phi_{\mathrm{A3}}(\bm r)
        =
        \int_0^1
        \psi_{\mathrm A}(\bm r;f)\,\mathrm df.
        \label{eq:ADensity}
    \end{equation}
    The B component is contributed only by the B block; therefore, $\phi_{\mathrm B}(\bm r)$ is obtained from Eq.~(\ref{eq:BlockDensity}) by setting $m=\mathrm B$. 


\PDFsubsection{Probability Density of the Sliding Degree of Freedom}
\label{sec:ProbabilityDensity}

    As follows from Eq.~(\ref{eq:ADensity}), the A-component density is obtained by integrating the topology-resolved density $\psi_{\mathrm A}(\bm r;f)$ over all accessible $f$.
    Since $\psi_{\mathrm A}(\bm r;f)$ is defined with the total single-chain partition function $Q=\int_0^1 Q_f\,\mathrm df$, the relative contributions of different $f$ are determined self-consistently by the equilibrium SCFT fields.
    Unlike a blend of fixed-topology copolymers, in which the overall populations of different topologies are prescribed by the blend composition, the distribution over $f$ in an L--S TT system is itself annealed and emerges self-consistently at equilibrium.


    
    Accordingly, for a finite spatial region $X$, we define the probability density of the sliding degree of freedom as
    \begin{equation}
        P_X(f)
        =
        \frac{
            \displaystyle
            \int_X
            \psi_{\mathrm A}(\bm r;f) \mathrm d\bm r
        }{
            \displaystyle
            \int_X
            \phi_{\mathrm A}(\bm r) \mathrm d\bm r
        },
        \label{eq:Probability}
    \end{equation}
    which satisfies $\int_0^1 P_X(f) \mathrm df=1$ according to Eq.~(\ref{eq:ADensity}).
    
    
    Here, $P_X(f)$ represents the probability density of chain topologies associated with A segments sampled within region $X$. 
    Because these A segments belong to chains characterized by different sliding configurations $f$, variations of $P_X(f)$ among different regions provide a spatially resolved statistical measure of how different topological configurations contribute to and redistribute within the ordered structure.
    
    The preferred topology associated with a region $X$ is defined as the location of the maximum of its topology distribution, $f_X^\ast=\arg\max_{f\in[0,1]}P_X(f)$.
    


    In this work, we consider three spatial regions $X\in\{\Lambda,\Omega(\phi_{\mathrm A}^{\mathrm{iso}}),\mathcal R(\bm r)\}$ to characterize the topology statistics at global, density-layer, and local spatial scales, respectively.
    The full-space region $\Lambda$ denotes the entire computational domain of the system.
    A finite-thickness density layer of the A component is defined as
    \begin{equation}
        \Omega(\phi_{\mathrm A}^{\mathrm{iso}})
        =
        \left\{
            \bm r\in\Lambda
            \,\middle|\,
            \left|\phi_{\mathrm A}(\bm r)-\phi_{\mathrm A}^{\mathrm{iso}}\right|
            <
            \Delta\,\phi_{\mathrm A}^{\mathrm{iso}}
        \right\}.
        \label{eq:IsodensityRegion}
    \end{equation}
    Here, $\phi_{\mathrm A}^{\mathrm{iso}}$ specifies the sampled A-component density level, and $\Delta$ specifies the relative half-width of the density layer.
    A local sampling region centered at position $\bm r$ is defined as
    \begin{equation}
        \mathcal R(\bm r)
        =
        \left\{
            \bm r'\in\Lambda
            \,\middle|\,
            \left|\bm r'-\bm r\right|
            \leq \Delta r
        \right\}.
        \label{eq:LocalRegion}
    \end{equation}
    The local topology distribution is understood in the limit $\Delta r\to0$.
    Provided that $\phi_{\mathrm A}(\bm r)>0$, this gives $P_{\mathcal R(\bm r)}(f)=\psi_{\mathrm A}(\bm r;f)/\phi_{\mathrm A}(\bm r)$, corresponding to the conditional topology distribution of A segments sampled at position $\bm r$.
    In numerical calculations, a small finite $\Delta r$ can be used as a spatial sampling tolerance to account for the discretization of the density field.
    The sliding degree of freedom $f$ itself remains an internal molecular variable rather than a spatial field.


\PDFsubsection{RPA Theory of L--S TT Copolymers}

    The RPA is a Gaussian fluctuation theory that describes phase stability through a second-order expansion of the free-energy functional.
    When applied around the homogeneous saddle point, it identifies the fluctuation mode that first destabilizes the disordered phase and determines the corresponding spinodal condition.\cite{Matsen1996Unifying, Wang2008Microphase, Li2020Largely}
    The comparison between the RPA spinodal and the SCFT order--disorder transition (ODT) provides a reference for assessing whether the transition behavior is mainly controlled by molecular correlations in the homogeneous state or by nonlinear effects in the ordered phase.\cite{Li2021Simplicity, Sides2004Continuous}
    For L--S TT copolymers, the topology-resolved RPA analysis additionally reveals the most unstable topology mode at the onset of instability, providing a microscopic reference for understanding the topology selection of ordered phases near the ODT.\cite{Zhao2024Phase}


    The saddle-point fields are denoted as $\{\phi_\alpha^\ast,\omega_\alpha^\ast\}$, and the density fluctuation is defined as $\delta\phi_\alpha(\bm r)=\phi_\alpha(\bm r)-\phi_\alpha^\ast$.
    Introducing these fluctuations in Fourier space gives the quadratic free-energy contribution:
    \begin{equation}
        \frac{\mathcal F^{(2)}}{N\rho_0}
        =
        \frac12
        \int
        \frac{\mathrm d\bm k}{(2\pi)^3}
        \delta\phi(\bm k)
        \Gamma_2(\bm k)
        \delta\phi(-\bm k),
        \label{eq:RPAKernel}
    \end{equation}
    where $\Gamma_2(\bm k)$ is the quadratic fluctuation kernel.
    For an incompressible two-component system, it can be expressed as $\Gamma_2(\bm k)=\widetilde C^{-1}(\bm k)-\chi N/2$, where $\widetilde C^{-1}(\bm k)$ originates from the single-chain correlation functions and represents the conformational entropy contribution, while $\chi N/2$ represents the enthalpic driving force for A--B segregation.
    Therefore, the RPA instability is determined by the competition between the conformational entropy encoded in the chain correlation functions and the enthalpic driving force for phase separation.


    For L--S TT copolymers, the single-chain correlations depend on the sliding degree of freedom.
    The component correlation functions entering the RPA kernel are obtained by averaging the fixed-topology correlations over the accessible topology ensemble:
    \begin{equation}
        \widetilde C_{\alpha\beta}(\bm k)
        =
        \int_0^1
        P_{\mathrm{RPA}}(f)
        \widetilde C_{\alpha\beta}^{\mathrm{Fix}}(\bm k,f)
        \mathrm df,
        \label{eq:RPACorrelation}
    \end{equation}
    where $\widetilde C_{\alpha\beta}^{\mathrm{Fix}}(\bm k,f)$ denotes the correlation function of a fixed topology $f$, and $P_{\mathrm{RPA}}(f)=Q_f[\omega_\alpha^\ast]/Q[\omega_\alpha^\ast]$ is determined by the saddle-point fields and describes the statistical weight of different sliding configurations.
    The inverse fluctuation kernel $\widetilde C^{-1}(\bm k)$ is subsequently constructed from these component correlations, and its detailed expression is provided in Supporting Information Sec.~\ntxref{SI-sec:RPA}.
    Therefore, the choice of $P_{\mathrm{RPA}}(f)$ determines how the sliding degree of freedom modifies the single-chain correlations and consequently the conformational entropy contribution to the fluctuation free energy.


    When RPA is used to determine the spinodal of the freely sliding system, the expansion is performed around the homogeneous saddle point.
    Because the homogeneous fields do not distinguish different sliding positions and changing $f$ only redistributes the connectivity among the A1, A2, and A3 blocks without changing the total A and B compositions, $Q_f[\omega_\alpha^\ast]$ is independent of $f$, giving $P_{\mathrm{RPA}}(f)=1$ over $[0,1]$.
    We denote this formulation as RPA--Slide, which describes the actual freely sliding L--S TT system, with the corresponding fluctuation kernel $\widetilde C_{\mathrm{Slide}}^{-1}(\bm k)$ constructed from the averaged component correlations.
    The unstable fluctuation wavelength is determined by
    \begin{equation}
        \left.
        \frac{\partial
        \widetilde C_{\mathrm{Slide}}^{-1}(k)}
        {\partial k}
        \right|_{k=k_{\mathrm{RPA\text{-}Slide}}^\ast}
        =
        0.
    \end{equation}
    The corresponding spinodal value is obtained from the instability condition $\widetilde C_{\mathrm{Slide}}^{-1}(k_{\mathrm{RPA\text{-}Slide}}^\ast)-N\chi^{\text{Slide}}_{\mathrm{sp}}/2=0$.


    For comparison, our previous treatment~\cite{Zhao2024Phase} employed the most-unstable-topology approximation, denoted as RPA--MUT, to identify the fixed-topology fluctuation mode that becomes unstable first and is expected to dominate the initial formation of ordered phases near the ODT.
    Instead of retaining the complete topology ensemble, this approximation replaces $P_{\mathrm{RPA}}(f)$ by $\delta(f-f_0)$ and evaluates each fixed-topology fluctuation channel separately.
    The resulting kernel $\widetilde C_{\mathrm{MUT}}^{-1}(k,f_0)$ therefore depends on the candidate topology $f_0$, and the first unstable topology is identified by
    \begin{equation}
        (k_{\mathrm{RPA\text{-}MUT}}^\ast,
        f_{\mathrm{RPA\text{-}MUT}}^\ast)
        =
        \operatorname*{arg\,min}_{k,f_0\in[0,1]}
        \widetilde C_{\mathrm{MUT}}^{-1}(k,f_0).
    \end{equation}
    Here, the optimization over $f_0$ searches among all fixed-topology fluctuation channels because the topology degree of freedom is not averaged as in RPA--Slide.
    
